% from ega4-2-fr.tex
\oldpage[IV-2]{80}

\emph{totale de $x$} (ou du
cycle premier $\overline{\{x\}}$) pour $\mathscr F$ (relative à $k$) le
produit de la multiplicité radicielle de $x$ pour $\mathscr F$ par la
multiplicité séparable de $k(x)$ sur $k$.
\end{definition}

Il revient au même de dire que la multiplicité totale de $x$ pour $\mathscr F$
est le produit de la longueur
$\operatorname{long}_{\mathcal O_x}(\mathscr F_x)$ par la multiplicité totale
de $k(x)$ sur $k$ (4.7.4). Avec les notations de (4.7.11), la multiplicité
totale de $x$ pour $\mathscr F$ est égale à la somme des multiplicités totales
pour $\mathscr F'$ des points maximaux $x'_i$ de
$\operatorname{Supp}(\mathscr F')$ qui sont au-dessus de $x$; en outre, il
existe une extension finie $k'$ de $k$, telle que la multiplicité totale de
$x$ pour $\mathscr F$ soit égale à
$\sum_i\operatorname{long}(\mathscr F'_{x'_i})$.

\begin{proposition}[4.7.13]
\label{IV.2.4.7.13-fr}
Les hypothèses et les notations étant celles de (4.7.10), pour que
$\mathscr F$ soit géométriquement ponctuellement intègre au point $x$, il
suffit que $x$ n'appartienne qu'à un seul cycle premier associé
(nécessairement non immergé) de $\mathscr F$ et que la multiplicité totale
pour $\mathscr F$ du point générique $z$ de ce cycle soit égale à $1$.
\end{proposition}

\begin{proof}
En effet, si $k'$ est une extension finie de $k$, $x'$ un point de
$X'=X\otimes_k k'$ au-dessus de $x$ et
$\mathscr F'=\mathscr F\otimes_k k'$, $x'$ ne peut appartenir qu'à un seul
cycle premier associé à $\mathscr F'$, puisqu'il n'y a qu'un seul point
maximal de $\operatorname{Supp}(\mathscr F')$ au-dessus de $z$ par hypothèse;
en outre, $\mathscr F$ est géométriquement réduit d'après (4.7.10), d'où la
conclusion.
\end{proof}
\subsection{Corps de définition}
\label{subsection:IV.4.8-fr}

\begin{env}[4.8.1]
\label{IV.2.4.8.1-fr}
Étant donné un préschéma $X$, nous désignerons, dans ce numéro, par
$\mathfrak S(X)$ l'ensemble des sous-préschémas de $X$, et comme d'ordinaire
par $\mathfrak P(X)$ l'ensemble des parties de l'espace sous-jacent à $X$;
pour tout $\mathcal O_X$-Module quasi-cohérent $\mathscr F$, nous désignerons
par $\Phi(\mathscr F)$ l'ensemble des sous-$\mathcal O_X$-Modules
quasi-cohérents de $\mathscr F$.
\end{env}

\begin{env}[4.8.2]
\label{IV.2.4.8.2-fr}
Soient $k$ un corps, $X$, $Y$ deux $k$-préschémas, $\mathscr F$, $\mathscr G$
deux $\mathcal O_X$-Modules quasi-cohérents, $K$, $K'$ deux extensions de $k$
telles que $K'\subset K$. On a alors, correspondant au morphisme
$\operatorname{Spec}(K)\to\operatorname{Spec}(K')$, des applications
canoniques
\begin{align}
\Phi(\mathscr F\otimes_kK')
&\longrightarrow\Phi(\mathscr F\otimes_kK),
\tag{4.8.2.1}
\label{IV.2.4.8.2.1-fr}\\
\operatorname{Hom}(\mathscr F\otimes_kK',\mathscr G\otimes_kK')
&\longrightarrow
\operatorname{Hom}(\mathscr F\otimes_kK,\mathscr G\otimes_kK),
\tag{4.8.2.2}
\label{IV.2.4.8.2.2-fr}\\
\mathfrak S(X_{(K')})&\longrightarrow\mathfrak S(X_{(K)}),
\tag{4.8.2.3}
\label{IV.2.4.8.2.3-fr}\\
\operatorname{Hom}_{K'}(X_{(K')},Y_{(K')})
&\longrightarrow\operatorname{Hom}_K(X_{(K)},Y_{(K)}),
\tag{4.8.2.4}
\label{IV.2.4.8.2.4-fr}\\
\mathfrak P(X_{(K')})&\longrightarrow\mathfrak P(X_{(K)}).
\tag{4.8.2.5}
\label{IV.2.4.8.2.5-fr}
\end{align}
Si $p_X:X_{(K)}\to X_{(K')}$ est la projection canonique, l'application
(4.8.2.5) est $M\mapsto p_X^{-1}(M)$, et de même (4.8.2.3) est
$Z\mapsto p_X^{-1}(Z)=Z\otimes_{K'}K$ (I, 4.4.1). L'application (4.8.2.1)
est
\[
\mathscr H\longmapsto p_X^*(\mathscr H)
=\mathscr H\otimes_{\mathcal O_{X_{(K')}}}\mathcal O_{X_{(K)}};
\]
(4.8.2.2) est $u\mapsto p_X^*(u)$ et enfin (4.8.2.4) est l'application
$f\mapsto f_{(K)}$.

Si l'on désigne par $\varepsilon_{K,K'}$ une quelconque de ces applications
canoniques, il est clair que si $K''$, $K'$, $K$ sont trois extensions de $k$
telles que $K''\subset K'\subset K$, on a
\begin{equation}
\varepsilon_{K,K''}=\varepsilon_{K,K'}\circ\varepsilon_{K',K''}.
\tag{4.8.2.6}
\label{IV.2.4.8.2.6-fr}
\end{equation}
\end{env}

